\documentclass[draft]{rureport}
\svnid{$Id: example-concept-generation.tex 133 2016-08-15 13:35:23Z foley $}
% use svn properties to set svn:keywords to for Id and URL (or HeadURL)

%% Fixmes are reminders of things that still need to be done.
%% The default fixmes are:  \fxnote{} \fxwarning{} \fxerror{} \fxfatal{} (same as \fixme{})
% if you want personalized fixmes, then register the authors here
% notice that the first field is 2 letters, the second is 3.
\FXRegisterAuthor{jf}{jtf}{foley}
% this registers \jfnote{}, \jfwarning{}, \jferror{}, \jffatal{}

%% declare the paths(s) where you graphics files can be found
\graphicspath{{graphics/}{Graphics/}{./}}

\author{Firstname13~Lastname\formatemail{email11}\\
  Firstname2~Lastname,firstname\formatemail{email12}}
\title{Concept Generation Template}  % The title, for the titlepage
\course{T-411-MECH Mechatronics 1}

\usepackage{hyperref} % must be last package loaded!
\begin{document}
\maketitle{}  % make a title page from the Title, Date, and Author
\listoffixmes{}
\subsection*{Instructions for Concept Generation}
This proposal is a team assignment.
Your job is to write up three ideas for a final project.
Each idea should contain a quick description, sketch or picture, and a link to a related article or product.
This document is provided as a template --- you will need to change it and comment out sections that are irrelevant.
The pictures provided here are just to demonstrate how graphics are included.
Your description should say what the item is, how it would be used, and what parts you think you might need.

For your own sanity while debugging, using SVN, and using a Grammarly grammar check, make sure you put only one sentence per line.
This means pressing ``enter'' on your keyboard every time you end a sentence.
SVN checks for conflicts and differences on a line-by-line basis.
\LaTeX{} shows error on a line-by-line basis.
If you make an entire paragraph be on a single line, you will have trouble finding any error on that line.

\subsection*{Using References/Citations}
List your references (journal articles, textbook pages, etc.).
A citation is probably not a good one if you can't at least get a title, year, and an author.
Publisher and date is important, but for online resources, this can be hard to find.
If you are using the standard template, this information should go into a \url{references.bib} file.
If you want to use a different file, you should change the \begin{verbatim}\bibliography{references}\end{verbatim} to have the name of the \path{.bib} file you chose.

Wikipedia and google are great resources for finding useful citations and pictures.
They are not good things to cite directly: anyone can edit any entry with a minimal level of checking.
Wikipedia is best used to find related articles that are the ones that you should be putting in your document.

We will use the IEEE citation style, referring to each reference with a consecutive number in the body of the text and listing the references in the same numerical order in the References section.
Here is a link to a website describing the IEEE citation style (from the University of Canterbury):

\url{http://library.canterbury.ac.nz/services/ref/ieee.shtml}

You can change the BibTeX style to IEEE by setting the bibliography style to
IEEEtran:\begin{verbatim}\bibliographystyle{IEEEtran}\end{verbatim}

As an example of a citation in text, Carryer et al.~\cite{carryer2011IntroMechatronics} is the textbook for the class.
This is another citation~\cite{kreher2005pseudocode}.

Dustcloud citation \cite{weber2012eyjafjallajokull} Sverrir says Thanks!!
\section{Killer Robot}
\begin{figure}
\centering
\includegraphics[height=50mm]{bender}
\caption{Concept art for Killer Robot prototype\label{fig:killerrobot}\cite{hall_virtuousmachines}}
\end{figure}

The killer robot is designed to enslave humanity.
A prototype can be
seen in Figure~\ref{fig:killerrobot}. Warbrick etal.~\cite{Warbrick94essentiallatex} contains LaTeX code for making such a system possible.
We would need 2 particle projection cannons, a nuclear reactor, and a pineal gland in order to make it work.

\section{Sharks with Lasers}
\begin{figure}
\centering
\includegraphics[height=50mm]{sharks_w_laser_beams}
\caption{Example of a finished Shark Laser\label{fig:sharklaser}\cite{ramirez2010sharkweek}}
\end{figure}
As mentioned in the movie ``Austin Powers'' we need some sharks with lasers!
These sharks would find our targets and eliminate them.
An Arduino based control and guidance system completed system is shown in Figure~\ref{fig:sharklaser}.
The wiring instructions are inspired by~\cite{carryer2011IntroMechatronics}.
We will need a swimming pool, two sharks, two arduinos, and two turbolasers to make this possible.

\section{Fake Moon landing}
\begin{figure}
\centering
\includegraphics[height=50mm]{moonhoax}
\caption{Design of a Moon Landing Hoax set\label{fig:fakemoonlanding}\cite{krystek2003moonhoax}}
\end{figure}

As everyone knows, the Moon landing in the 1960's was faked.~\cite{krystek2003moonhoax}
Here's how it was done.
A prototype green-screen was used with custom props and rigging to emulate the lunar environment.
You can see the original implementation in Figure~\ref{fig:fakemoonlanding}.
This design could be used to fake other events, such as the war in Iraq.
We would need a video frame grabber and an Arduino to build this project.


\bibliographystyle{IEEEtran}
\bibliography{references}

\end{document}
%%%%%%%%%%%%%%%%%%%% TeXStudio Magic Comments %%%%%%%%%%%%%%%%%%%%%
%% These comments that start with "!TeX" modify the way TeXStudio works
%% For details see http://texstudio.sourceforge.net/manual/current/usermanual_en.html   Section 4.10
%%
%% What encoding is the file in?
% !TeX encoding = UTF-8
%% What language should it be spellchecked?
% !TeX spellcheck = en_US
%% What program should I compile this document with?
% !TeX program = pdflatex
%% Which program should be used for generating the bibliography?
% !TeX TXS-program:bibliography = txs:///bibtex
%% This also sets the bibliography program for TeXShop and TeXWorks
% !BIB program = bibtex

%%% Local Variables:
%%% mode: latex
%%% TeX-master: t
%%% End:
